% =====================================================================
% Table 3 --- the 5 x 3 matrix: the closure conditions (DRAFT)
% ---------------------------------------------------------------------
% Extends Table~\ref{tab:closure-conditions} from the two pole closures
% (2 columns) to the full five labour closures (5 rows) crossed with the
% three financing closures (3 columns), as announced in Section 3
% ("Before we map these five closures in a monumental table ...").
% Structure mimics Table 2: each cell carries the two equations that
% define that economy -- one from the short-run (labour) block, one from
% the financing block. The labour condition is the row, the financing
% condition is the column; the financing equation therefore repeats down
% a column, exactly as the short-run equation repeated down its column in
% Table 2.
%
% Ordering of rows follows the registry id order (ALPHA, BETA, GAMMA,
% DELTA, BF) = standard CGE, elastic (sectoral) supply, employment adjustment,
% Leontief, no-reallocation. Provisional -- reorder if the narrative
% prefers the two poles first.
%
% NAMING (to match the pending naming sweep): the row labels use the
% conventional vocabulary -- "standard CGE closure", "elastic labour
% supply", "employment adjustment closure", "Leontief closure",
% "no-reallocation closure".
% =====================================================================
\begin{table}[htbp]
\centering
\footnotesize
\setlength{\tabcolsep}{3pt}
\renewcommand{\arraystretch}{1.35}
\begin{tabular}{@{}>{\raggedright\arraybackslash}p{2.30cm}
                  >{\raggedright\arraybackslash}p{3.90cm}
                  >{\raggedright\arraybackslash}p{3.90cm}
                  >{\raggedright\arraybackslash}p{3.90cm}@{}}
\toprule
\textbf{Labor-market closure}
& \textbf{Financing: F1}\newline \scriptsize repurposing of consumption
& \textbf{Financing: F2}\newline \scriptsize refinancing via taxes
& \textbf{Financing: F3}\newline \scriptsize external finance \\
\midrule

\textbf{Standard CGE}\newline
  \scriptsize full employment, one flexible wage, mobile labor ($\eta = 1$)
& $L_i = L_i^{\mathrm{cm}}(p, y_i, w)$,\quad $\sum_i L_i = \bar L$
  \newline $\sum_i p_i c_i^{h} = E^{h}$
& $L_i = L_i^{\mathrm{cm}}(p, y_i, w)$,\quad $\sum_i L_i = \bar L$
  \newline $T(p) = \sum_i p_i g_i$
& $L_i = L_i^{\mathrm{cm}}(p, y_i, w)$,\quad $\sum_i L_i = \bar L$
  \newline $B_{\mathrm{gov}} = \sum_i p_i g_i$ \\
\midrule

\textbf{Elastic labor supply}\newline
  \scriptsize behavioral supply on the real wage; $N$ sectoral markets
& $L_i = \bar L_i\,[\,(w_i/\mathcal{P}(p))/(w_i^{0}/\mathcal{P}^{0})\,]^{\eta_s}$
  \newline $\sum_i p_i c_i^{h} = E^{h}$
& $L_i = \bar L_i\,[\,(w_i/\mathcal{P}(p))/(w_i^{0}/\mathcal{P}^{0})\,]^{\eta_s}$
  \newline $T(p) = \sum_i p_i g_i$
& $L_i = \bar L_i\,[\,(w_i/\mathcal{P}(p))/(w_i^{0}/\mathcal{P}^{0})\,]^{\eta_s}$
  \newline $B_{\mathrm{gov}} = \sum_i p_i g_i$ \\
\midrule

\textbf{Employment adjustment}\newline
  \scriptsize wage pinned, employment uncapped
& $w = 1$ (numeraire),\quad $L = \sum_i L_i$
  \newline $\sum_i p_i c_i^{h} = E^{h}$
& $w = 1$ (numeraire),\quad $L = \sum_i L_i$
  \newline $T(p) = \sum_i p_i g_i$
& $w = 1$ (numeraire),\quad $L = \sum_i L_i$
  \newline $B_{\mathrm{gov}} = \sum_i p_i g_i$ \\
\midrule

\textbf{Leontief closure}\newline
  \scriptsize wage pinned and Leontief limit, quantities absorb
& $w = 1$,\quad $(\theta, \epsilon, \sigma) \to 0^{+}$
  \newline $\sum_i p_i c_i^{h} = E^{h}$
& $w = 1$,\quad $(\theta, \epsilon, \sigma) \to 0^{+}$
  \newline $T(p) = \sum_i p_i g_i$
& $w = 1$,\quad $(\theta, \epsilon, \sigma) \to 0^{+}$
  \newline $B_{\mathrm{gov}} = \sum_i p_i g_i$ \\
\midrule

\textbf{No-reallocation}\newline
  \scriptsize immobile labor, sectoral wages free ($\eta = 0$)
& $\eta = 0$:\quad $L_i = \bar L_i$, $w_i$ free
  \newline $\sum_i p_i c_i^{h} = E^{h}$
& $\eta = 0$:\quad $L_i = \bar L_i$, $w_i$ free
  \newline $T(p) = \sum_i p_i g_i$
& $\eta = 0$:\quad $L_i = \bar L_i$, $w_i$ free
  \newline $B_{\mathrm{gov}} = \sum_i p_i g_i$ \\
\bottomrule
\end{tabular}
\caption{The full closure matrix: the five labor-market closures (rows)
crossed with the three financing closures (columns). Each cell states the
two closure conditions that define that economy --- one for the short run
(labor market, the row), one for the financing of the program (the
column). The standard CGE row and the $F1$/$F2$ column reproduce
Table~\ref{tab:closure-conditions}; the remaining rows add the labor-market
variations between the two poles --- behavioral supply on the real wage
(elastic labor supply), an uncapped extensive margin (employment adjustment),
the fixed-coefficient limit (Leontief), and immobile labor with sectoral
wages (no-reallocation). Notation follows Section~\ref{sec:kernel};
$\mathcal{P}(p)$ is the consumption price index, $E^{h}$ household consumption
expenditure, $T(p)$ the program tax, and $B_{\mathrm{gov}}$ the externally
booked program cost.}
\label{tab:closure-matrix}
\end{table}